\section{The expected number of pairs at large hopcount}\label{sec:profile}

Fix a root $r\in[n]$. The union of the minimum-weight paths from $r$ to the other vertices is a spanning tree, the
shortest-path tree of $r$. Explore it by increasing distance from $r$. When $m$ vertices have been found, the
memoryless property of the exponential distribution makes the residual weights of the $m(n-m)$ edges between found
and unfound vertices, that is, the weight of each such edge less the time elapsed since its found endpoint was
reached, independent exponential variables of mean one, so the next vertex is attached by an edge chosen
uniformly among them, and its parent is uniform among the $m$ vertices found so far, independently of the past. The
tree, with its vertices numbered in the order in which they are found, is therefore a random recursive tree: vertex
$m+1$ is attached to a uniform vertex among $1,\dots,m$~\cite{vdHHVM,SM}. The hopcount $H_n(r,v)$ is the depth of $v$
in this tree. The trees of different roots are dependent, and nothing below uses more than the law of one of them.

Let $Z_{n,d}$ be the number of vertices at depth $d$ in a random recursive tree with $n$ vertices, and
\[
F_n(z)=\E\sum_{d\ge0}Z_{n,d}\,z^d,\qquad f_{n,d}=[z^d]F_n(z)=\E Z_{n,d}.
\]
A new vertex attached to a uniform vertex among $n$ has depth one more than its parent, so
$F_{n+1}(z)=(1+z/n)F_n(z)$ and
\begin{equation}\label{eq:F}
F_n(z)=\prod_{j=1}^{n-1}\Bigl(1+\frac zj\Bigr)=\frac{\Gamma(n+z)}{\Gamma(n)\,\Gamma(1+z)} .
\end{equation}
Thus $(n-1)!\,f_{n,d}$ is the unsigned Stirling number of the first kind $\genfrac{[}{]}{0pt}{}{n}{d+1}$; this description of the expected profile is
classical, see~\cite{SM}, \cite[Chapter~6]{Drmota} and~\cite{FHN}.

\begin{lemma}\label{lem:first}
For all $n\ge2$ and $k\ge1$,
\[
\E N_{n,k}=\frac n2\sum_{d\ge k}f_{n,d}.
\]
\end{lemma}

\begin{proof}
Each pair $\{i,j\}$ is counted once from the root $i$ and once from the root $j$, so
$2N_{n,k}=\sum_{r}\#\{v\ne r:H_n(r,v)\ge k\}$. For each $r$ the inner number is the number of vertices at depth at
least $k$ in the shortest-path tree of $r$, whose mean is $\sum_{d\ge k}f_{n,d}$.
\end{proof}

The sum in Lemma~\ref{lem:first} is over the whole profile. It should not be replaced by $n-1$ times the tail of the
hopcount between two fixed vertices computed from $F_n(z)/n$: that function is the generating function of the depth
of a uniform vertex with the root included, and the depth of a uniform vertex other than the root has generating
function $(F_n(z)-1)/(n-1)$.

\subsection*{Tilting}
Let $a>0$. Let $B_1,\dots,B_{n-1}$ be independent Bernoulli variables with
\[
\Prob(B_j=1)=q_j=\frac{a}{j+a},\qquad S_n=\sum_{j=1}^{n-1}B_j .
\]
Since $(1+au/j)/(1+a/j)=1-q_j+q_ju$, we have $\E u^{S_n}=F_n(au)/F_n(a)$ and therefore
\begin{equation}\label{eq:tilt}
f_{n,d}=F_n(a)\,a^{-d}\,\Prob(S_n=d)\qquad(d\ge0).
\end{equation}
For $a=1$ this is the representation of the depth of a uniform vertex of the random recursive tree as a sum of
independent indicators~\cite{Devroye88}.
The mean and variance of $S_n$ are
\[
\mu_n=\sum_{j<n}\frac{a}{j+a},\qquad V_n=\sum_{j<n}\frac{aj}{(j+a)^2} .
\]
From $0\le a/j-a/(j+a)\le a^2/j^2$ and $|aj/(j+a)^2-a/j|\le 2a^2/j^2+a^3/j^3$,
\begin{equation}\label{eq:muV}
\mu_n=a\log n+O(1),\qquad V_n=a\log n+O(1),
\end{equation}
with constants depending on $a$ only. As $V_n>0$ for every $n\ge2$, there is $c_a>0$ with $V_n\ge c_a\log n$ for all
$n\ge2$.

\begin{lemma}\label{lem:atom}
Let $X$ be a sum of finitely many independent Bernoulli variables, with variance $V>0$. Then for all real $t$,
\begin{equation}\label{eq:cf}
\bigl|\E e^{itX}\bigr|\le\exp\bigl(-2V\sin^2(t/2)\bigr),
\end{equation}
and $\sup_{d}\Prob(X=d)\le\sqrt{\pi/(8V)}$. In particular, for all $n\ge2$,
\begin{equation}\label{eq:atom}
\sup_{d}\Prob(S_n=d)\le\sqrt{\frac{\pi}{8V_n}}\le\frac{C_a}{\sqrt{\log n}} .
\end{equation}
\end{lemma}

\begin{proof}
For one Bernoulli variable with parameter $q$, $|1-q+qe^{it}|^2=1-4q(1-q)\sin^2(t/2)$, and
$\sqrt{1-y}\le e^{-y/2}$ for $0\le y\le1$; multiplying over the summands gives~\eqref{eq:cf}. By Fourier inversion
for an integer-valued variable,
\[
\Prob(X=d)=\frac1{2\pi}\int_{-\pi}^{\pi}\E e^{itX}\,e^{-itd}\,dt
\le\frac1{2\pi}\int_{-\pi}^{\pi}\bigl|\E e^{itX}\bigr|\,dt
\le\frac1{2\pi}\int_{-\infty}^{\infty}e^{-2Vt^2/\pi^2}\,dt=\sqrt{\frac{\pi}{8V}},
\]
where we used $\sin(|t|/2)\ge|t|/\pi$ for $|t|\le\pi$.
\end{proof}
